\subsection{The problem and its development}

Brennan's 1978 work \cite{Brennan1978,Brennan1978LOMI} improved the
previously known upper range beyond $s=3$, as recorded in
Bertilsson's historical account \cite[Chapter~3, pp.~43--44]{Bertilsson1999}.
The standard formulation asks for the full interval $4/3<s<4$ on every
simply connected plane domain admitting a conformal bijection to the disk.
Neither boundedness nor boundary regularity is imposed. The classical
slit-plane obstructions at both endpoints are verified in
Section~\ref{sec:integrability}.

Subsequent progress connected the area question to finer estimates for
derivative integral means. Pommerenke developed this approach in
\cite{Pommerenke1985IntegralMeans,Pommerenke1985IntegralMeansII}.
Carleson and Makarov related negative integral means to the distribution
of large harmonic measure and proved the linear spectral tail
$B_{\SSS}(t)=|t|-1$ for all sufficiently negative $t$
\cite[\S1.4]{CarlesonMakarov1994}. The proposed threshold $t=-2$
is the one associated with Brennan's conjecture
\cite{HedenmalmShimorin2005}. The statement for sufficiently negative parameters alone
does not reach that endpoint.

Bertilsson studied coefficient estimates for negative powers of
derivatives of univalent functions \cite{Bertilsson1998}. His thesis
reports the area range $4/3<s<3.421$ \cite{Bertilsson1999} and
establishes equivalences between the spectral endpoint and
uniform integral-means formulations
\cite[Theorem~3.7, especially parts (a) and (f), pp.~44--45]{Bertilsson1999}.
Thus uniform estimates have an established place in the formulation of
Brennan's problem. In the present argument, the common constant in
\eqref{eq:main-uniform} arises directly from an operator on the entire
compact normalized class.

Hedenmalm and Shimorin combined area inequalities with weighted Bergman
space methods to improve general bounds for the integral-means spectrum.
Their published table gives $B_{\SSS}(-2)\le1.218$
\cite[Table~1, p.~385]{HedenmalmShimorin2005}.
Sola extended the Bergman-space computation
\cite{Sola2006}, reporting the improved estimate
$B_{\SSS}(-2)\le1.206$ in his thesis \cite[p.~6]{Sola2007Thesis}.
These estimates illustrate the remaining gap between general analytic
bounds and the conjectured exponent $1$.

Recent works treat the conjectured integrability range under additional
hypotheses. Aleman and Kouroupis prove it for members of continuous
semigroups of conformal disk self-maps and for Koenigs maps of nontrivial
such semigroups
\cite[Theorem~1.1 and Proposition~2.2]{AlemanKouroupis2025}.
Zheng treats simply connected polynomial basins at infinity for
$P(z)=z^m+a_0$, $m\ge3$, and for monic degree-$m$ polynomials whose
coefficients of $z,\ldots,z^{m-1}$ have modulus smaller than a
degree-dependent threshold; these basins
are considered in the sphere and contain infinity
\cite[Theorems~1.3--1.4]{Zheng2026Basins}.
Jin computes the real integral-means spectra of normalized rational
functions univalent in the disk,
with inverse-square supremum $1$ over that class
\cite[Theorem~6.20 and Corollaries~6.22--6.23]{Jin2025IntegralMeans}.
His work on Schwarzian multiplication operators recalls Shimorin's
sufficient multiplier-norm criterion for the inverse-square endpoint
of an individual normalized univalent function
\cite[Proposition~1.4(2)]{Jin2026Multipliers}.
Theorem~\ref{thm:main} concerns the full class of univalent disk maps.

\subsection{Analytic and geometric ingredients}

Normalized changes of base point, distortion estimates, and compactness
of the schlicht class are classical tools of geometric function theory
\cite[Chapters~1--2]{Duren1983}. A related probabilistic method for
exterior integral means appears in the conformal snowflakes of Beliaev
and Smirnov. Their positive integral operator gives lower bounds for an
averaged exponent defined from expected derivative integral means.
For each real $t$ and $\varepsilon>0$, they construct a random snowflake
whose averaged exponent exceeds $B(t)-\varepsilon$, where $B$ is their
universal spectrum defined using exterior Riemann maps
\cite[Lemma~3.4 and Theorem~4.1]{BeliaevSmirnov2010}.
Here the operator acts on all normalized half-plane maps, and
hypothetical excess growth produces a weighted affine law. The
operator is defined in Section~\ref{sec:preliminaries};
Section~\ref{sec:weighted-law} constructs its eigenmeasure and proves
the averaging identity.

The critical-point comparison is a planar index argument. Its
boundary calculation uses Hopf's secant proof of the tangent-turning
theorem \cite[\S2, pp.~53--55]{Hopf1935}, while the regular-target
step is an elementary plane case of Sard's theorem
\cite[Theorem~4.1]{Sard1942}.
We keep the index of a possible pole separate from those of the finite
critical points. The pairing is defined by ranks in ordered lists;
its measurability is proved through continuous local inverse branches,
and its equivariance follows from the ordering itself.

The subsequent area comparison is related in form to equivariant
mass transport, including the general framework for nonunimodular
group actions developed by Gentner and Last \cite{GentnerLast2011}.
We use a direct calculation with local changes of variables and
explicit expanding boxes. The affine covariance, exceptional-target
transfer, and boundary-cost estimates needed here are all proved
within the paper.

An independent companion applies a related affine-law and critical-point
comparison to prove a uniform inverse-first-power bound with exponent
strictly below $1/4$ for normalized univalent disk maps, disproving
Kraetzer's bounded-spectrum prediction at $p=-1$
\cite[Theorem~1.1]{OpenAIStrictInverseFirst2026}.
The proof here does not use that strict estimate; its H\"older
interpolation from \eqref{eq:main-uniform} gives exponent $1/2$ with
an arbitrary positive loss at $t=-1$.
